%% kapa: exam=1 round=22 year=2011 period=1 form=A subject=civil_law subject_ko=민법 seq=33 number=33
%% source: Q-Net 감정평가사 기출문제 (cst00309) artlSeq=5036790 file=extracted/1008/5036790/2011년 제22회 감정평가사 1차시험 기출문제/1교시.hwp
%% conversion: kordoc 4.24.51 -> tools/gen1.py (auto); review=proofread-sonnet; flags=none
\begin{question}{33}
등기의 추정력에 관한 설명으로 옳지 않은 것은? (다툼이 있으면 판례에 의함)
\begin{choices}
  \item 부동산에 관한 소유권이전등기가 있는 이상 그 절차는 정당한 것으로 추정된다.
  \item 소유권이전의 등기명의자는 제3자뿐만 아니라 그 전소유자에 대하여도 적법한 등기원인에 의하여 소유권을 취득한 것으로 추정된다.
  \item 가등기에는 가등기원인에 대한 적법추정력이 인정되지 않는다.
  \item 근저당권의 설정등기가 되어 있으면 이에 상응하는 피담보채권의 존재가 추정된다.
  \item 소유권보존등기의 명의인이 부동산을 양수받은 것이라 주장하고 전소유자는 양도사실을 부인하는 경우에도 그 보존등기의 추정력은 인정된다.
\end{choices}
\end{question}
\answer{5}
